\subsection{Modely zaměřené na zdrojový kód}
\label{subsec:models-code}

První skupinu tvoří modely, které byly od počátku navrženy nebo dodatečně dotrénované (\textit{fine-tuned}) specificky pro práci se zdrojovým kódem.
Tyto modely mají oproti obecným jazykovým modelům výhodu v tom, že jejich trénovací data obsahují výrazně vyšší podíl zdrojového kódu, dokumentace k API a diskusí o programátorských problémech včetně potenciálních chyb.
Díky tomu lépe rozumí syntaktickým konstrukcím, standardním vzorům a typickým chybám v různých programovacích jazycích.

\begin{itemize}
    \item \textbf{Code Llama}~\cite{codellama} je rodina modelů od společnosti Meta, odvozená z obecného modelu LLaMA~2.
    Modely byly dále trénovány na 500 miliardách tokenů zdrojového kódu a jsou k dispozici ve velikostech 7B, 13B, 34B a 70B parametrů.
    Varianta \textit{Code Llama Instruct} je optimalizována pro sledování instrukcí, což je pro účely analýzy kódu na základě zadaného promptu obzvláště relevantní.
    Na benchmarku HumanEval~\cite{humaneval} dosahuje varianta 34B přibližně 48\,\% úspěšnosti (pass@1).

    \item \textbf{DeepSeek Coder}~\cite{deepseek-coder} je rodina modelů trénovaných od základu na 2 bilionech tokenů, z nichž 87\,\% tvoří zdrojový kód a zbývajících 13\,\% přirozený jazyk.
    Modely jsou dostupné ve velikostech od 1,3B do 33B parametrů.
    Na benchmarku HumanEval dosahuje varianta 33B Instruct přibližně 79\,\% úspěšnosti, což z ní činí jeden z nejvýkonnějších open-source modelů pro generování kódu.

    \item \textbf{StarCoder~2}~\cite{starcoder2} vznikl v rámci projektu BigCode jako plně otevřený model trénovaný na datové sadě The Stack v2~\cite{kocetkov-stack} o rozsahu 3,3 bilionu tokenů.
    Je k dispozici ve velikostech 3B, 7B a 15B parametrů.
    Největší varianta dosahuje na HumanEval přibližně 46\,\% úspěšnosti.
    Výhodou StarCoderu~2 je transparentnost trénovaných dat, takže vývojáři mohou ověřit, zda byl jejich kód použit pro trénování, případně požádat o jeho vyřazení.

    \item \textbf{Qwen2.5-Coder}~\cite{qwen25-coder} je model od společnosti Alibaba, který patří do rodiny Qwen a je specializovaný na práci s kódem.
    Byl trénován na 5,5 bilionu tokenů zdrojového kódu pokrývajících více než 90 programovacích jazyků.
    Model je dostupný ve velikostech od 0,5B do 32B parametrů a na HumanEval dosahuje varianta 32B Instruct přibližně 92\,\% úspěšnosti, čímž se řadí mezi nejlepší open-source modely v oblasti generování a analýzy kódu.
    Navazující verze \textbf{Qwen3-Coder}~\cite{qwen3-coder} přináší další vylepšení v oblasti porozumění kontextu a podpory pro více programovacích jazyků.
\end{itemize}

Společným znakem všech uvedených modelů je, že jsou dostupné jako open-source a lze je provozovat lokálně prostřednictvím nástrojů popsaných v \secref[sekci]{subsec:deployment-on-premise}.

\endinput